\section{Quantum interfaces: readout, unitary evolution, and bandwidth}
\label{sec:quantum_interfaces}

The effective gravitational equations of CAP-II live at the coarse-grained level.
Quantum mechanics enters as an interface theory for finite observers: the observer has an effective Hilbert space, applies local unitaries (scan/update), and accesses outcomes through finite-resolution readout instruments.
This section records the minimum interface structure used later, emphasizing what is assumed and what is merely a change of representation.

\subsection{Effective observer sector and POVM readout}
\label{subsec:povm_readout}

Let $\mathcal{H}_{\mathrm{eff}}$ be an effective observer sector and let $\rho$ be a density operator on $\mathcal{H}_{\mathrm{eff}}$.
Finite-resolution readout at scale $\varepsilon$ is modeled by a POVM $\{E_k^{(\varepsilon)}\}_{k}$ with $\sum_k E_k^{(\varepsilon)}=\ind$, giving Born probabilities
\begin{equation}
  P_k^{(\varepsilon)}=\Tr(\rho\,E_k^{(\varepsilon)}).
\end{equation}
The corresponding instrument (state update rule) may be written in Kraus form $\rho\mapsto M_k\rho M_k^\dagger/P_k$ with $E_k=M_k^\dagger M_k$.
Any such POVM can be implemented by a dilation (system--ancilla unitary followed by a projective measurement on the ancilla).
This is a standard consequence of Naimark and Stinespring representation theorems; see, e.g., \cite{Naimark1940,Stinespring1955,Kraus1983,NielsenChuang2000}.

\subsection{Bandwidth and controlled continuum operators}
\label{subsec:bandwidth_operators}

Assumption~\ref{ass:kernel_readout} asserts that readout is sufficiently regular and band-limited so that discrete difference operators approximate continuum derivatives on coarse-grained observables.
Operationally, this is the condition under which a continuum effective field theory is meaningful: the observer cannot resolve sub-$\varepsilon$ microstructure, and the induced operators act on $F_\varepsilon$ with controllable error when $h\ll\varepsilon$.
Appendix~\ref{app:readout_error_bounds} records a representative error estimate for gradients/Laplacians under band-limited kernel readout.

\subsection{Unitary evolution and representation choices}
\label{subsec:unitary_representation}

Whenever a one-parameter unitary family $U_t$ implements an automorphism of the effective algebra by $A\mapsto U_t^\dagger A U_t$, one may represent the same expectation data either by evolving observables (Heisenberg) or by evolving the state representative (Schr\"odinger).
CAP-II does not claim to derive this unitary structure from Einstein gravity; rather, it treats it as part of the observer-interface layer and uses it to define measurable quantities (e.g.\ scattering delay) that can be calibrated against the operational lapse.


